% Nouveau EGA V (cote 25): archive cote 25, site page 61 (PDF page 62).
% typescript. Draft compared: none.
% AI-transcribed by Claude Opus 5.5 (claude-opus-5-5, adaptive thinking, effort high), Message Batches API, batch ega25b (msgbatch_01A7n5YuXubUPNK1zsepi8RG), 2026-11-13; not yet checked by eye.
\scanpage{25/61}
\mspage{56}

L'énoncé \unsure{10.2.} est une conséquence à peu près immédiate du suivant,
qui est tout à fait indépendant d'histoires de sections hyperplanes, et
serait sans doute mieux à sa place dans \xout{xx} un nº en rabiot aux
``immersions régulières''.

\deleted{\uline{Proposition} \unsure{10.3.} Soient X un préschéma, \xout{$\underline{E}$} $\underline{G}$ un Module quasi-cohérent}
\deleted{sur X, et $u\colon \underline{G}\to\underline{O}_X$ un homomorphisme, $\underline{J} = u(G)$, $T=V(\underline{J})$. Soit}
\deleted{$\widetilde{X}$ déduit de X par éclatement de T. Considérons d'autre part}
\deleted{$P = P(\underline{G})$, l'homomorphisme canonique $\underline{G}_P\to\underline{O}_P(1)$ et son noyau}
\deleted{$\underline{H}$ ,(de sorte que nous avons la suite exacte $0\to\underline{H}-\underline{G}_P\to\underline{O}_P(1)\to 0$),}
\deleted{l'homomorphisme $u_P : \underline{G}_P \to \underline{O}_P$ , et l'Idéal qu coh $\underline{K} = u_P(\underline{H})\subset\underline{O}_P$ .}
\deleted{Alors $\widetilde{X}$ est canoniquement isomorphe à un sous-préschéma fermé de}
\deleted{$V(\underline{K})$ . Lorsque G est localement libre et u est ``\xout{re}régulier'', alors}
\deleted{l'isomorphisme précédent est un isomorphisme de $\widetilde{X}$ avec $V(\underline{K})$ lui-même,}
\deleted{de plus dans ce cas $\underline{H}$ est localement libre sur P et \xout{xxx} $\underline{H}\to\underline{O}_P$}
\deleted{(dont le\xout{\unsure{t}} préschéma d es zéros e st $\widetilde{X}$ ) est également régulier.}

\deleted{Le premier énoncé est à peu près trivial. Le deuxième est un}
\deleted{exercice qui ne peut offrir de difficulté (je ne l'ai pas fait en}
\deleted{détail, pensant que tu t'en tireras aussi bien que moi).}

Lorsque dans 10.2. \deleted{les fibres de X/S sont de} \inserted{$S = \mathrm{Spec}\,k$ et X est de} dimension $\unsure{\leqslant}1$,
alors l'hypothèse de \xout{\unsure{ex}}régularité faite équivaut à $T = \emptyset$ , de sorte
que $Y_{\underline{L}}\to X$ est un \uline{isomorphisme} . On trouve ainsi\unsure{;} en conjuguant
avec 9.2. :

\uline{Corollaire} \unsure{10.4.} (à \unsure{10.2.}) Soit X une courbe lisse géom connexe dans
\xout{dimensionxxx} l'espace projectif P sur un corps infini k, et soit $n\hand{\geqslant}2$.
Alors il existe un pinceau linéaire de n-formes sur P, définissant un
morphisme $X\to P^1$ ayant les propriétés suivantes: ce morphisme est

\mspage{61}
